Retrieval-Augmented Generation (RAG) is the family of architectures in which a large language model (LLM) is conditioned at inference time on a small set of documents drawn from an external corpus by a learned or hand-crafted retriever. The phrase was coined by Patrick Lewis and colleagues at Facebook AI Research in their NeurIPS 2020 paper \emph{``Retrieval-Augmented Generation for Knowledge-Intensive NLP Tasks''}, which introduced the canonical pipeline of a Dense Passage Retriever (DPR) feeding a BART-large generator and proposed two marginalization variants: RAG-Sequence, in which the same retrieved document is reused for every token of the answer, and RAG-Token, in which a different document may be used at each decoding step. The motivation for the design was empirical: pure parametric models such as T5 and BART memorize a fixed cross-section of world knowledge in their weights, but they cannot be updated, audited, or compressed after pretraining without expensive retraining. By contrast, a non-parametric memory of textual passages can be edited, version-controlled, and grown without modifying the model parameters, while the parametric LLM remains responsible for fluency, reasoning, and surface generation. This split is now a standard reference point in surveys by Yunfan Gao et al.~(2023, arXiv:2312.10997), Wenqi Fan et al.~(KDD 2024), Yizheng Huang and Jimmy Huang (ACM Computing Surveys 2026), and the trustworthy-RAG survey by Bo Ni et al.~(arXiv 2502.06872, 2025).
